\section{Reprodução dos experimentos}\label{ap:reproducao}

O repositório do projeto contém o código, os dados brutos e os scripts
necessários para reproduzir todos os resultados. Com Python~3.12 e Docker ou
Podman, a sequência é:

\begin{small}
\begin{verbatim}
uv venv --python 3.12 .venv && source .venv/bin/activate
uv pip install -e ".[dev]"            # versões fixadas em requirements-lock.txt
python -m pytest -q                   # testes automatizados
scripts/broker.sh start               # Mosquitto 2.0.20 em 127.0.0.1:1883
python -m experiments.run E1 E2 E3 E4 --workers 5
python -m experiments.run E5 --workers 3
python -m experiments.mqtt_integration
python -m experiments.replay && python -m experiments.analysis
python -m experiments.figures && python -m experiments.animations
python scripts/verify_reproducibility.py
scripts/build_paper.sh                # TeX Live 2026 em contêiner
\end{verbatim}
\end{small}

O executor grava uma linha por execução em \texttt{experiments/raw/E*.jsonl},
com a configuração completa, a semente e as métricas, e pula execuções já
concluídas, o que permite interromper e retomar a rodada. Os manifestos
registram o \emph{commit} do código, as versões das bibliotecas e a
plataforma.

\section{Material suplementar}\label{ap:suplementar}

As animações abaixo foram geradas a partir de execuções gravadas e estão no
diretório \texttt{animations/} do repositório, em GIF e MP4. Nenhuma delas é
ilustração conceitual: os quadros reproduzem os estados registrados.

\begin{description}
  \item[A\_enchimento] Ocupação de seis contêineres ao longo de 7 dias
  (política A, semente 1001).
  \item[B\_telemetria\_mqtt] Telemetria recebida pela plataforma pelo broker
  Mosquitto em uma execução do modo \emph{live} (80 contêineres virtuais), com
  eventos de limiar.
  \item[C\_coleta\_veiculos] Veículos executando as rotas da política D na
  malha de Aracaju.
  \item[D\_comparacao\_A\_D] Comparação sincronizada das políticas A e D com a
  mesma demanda.
  \item[E\_falha\_comunicacao] Estado verdadeiro e estado visto pela plataforma
  durante uma interrupção regional de 48~h.
\end{description}

\section{Mapa de pinos do firmware 2.0}\label{ap:pinos}

\begin{table}[ht]
\centering
\caption{Ligações do protótipo com o firmware 2.0 (DOIT ESP32 DevKit V1).}
\footnotesize
\begin{tabular}{@{}lll@{}}
\toprule
Sinal & GPIO & Observação \\
\midrule
HC-SR04 TRIG & 14 & nível alto de 3,3~V é suficiente para o disparo \\
HC-SR04 ECHO & 27 & via divisor 1~k$\Omega$ (série) / 2~k$\Omega$ (terra): $5 \times 2/3 \approx 3{,}33$~V \\
DHT11 DATA & 4 & alimentação em 3,3~V; \emph{pull-up} de 4,7--10~k$\Omega$ para 3,3~V \\
LED & 2 & LED da placa; saída \\
\bottomrule
\end{tabular}
\end{table}

\section{Declaração de uso de ferramentas de inteligência artificial}\label{ap:ia}

O desenvolvimento desta versão contou com a assistência de um modelo de
linguagem (Claude, da Anthropic) na revisão bibliográfica (busca e conferência
de metadados), na escrita e revisão do código, na execução dos experimentos e
na redação do texto. Todas as referências tiveram os metadados conferidos em
fontes oficiais (Crossref, OpenAlex ou página do editor), e o registro dessa
verificação, incluindo quais textos foram lidos integralmente e quais apenas
pelo resumo, está em \texttt{paper/references/biblioteca.csv}. Os resultados
numéricos foram produzidos pelos programas do repositório e inseridos no texto
automaticamente. A responsabilidade pelo conteúdo é dos autores, e o uso dessa
assistência deve observar as normas da instituição.
